
Our retrieval agent is built on top of ArtifactLinker \citep{yu2026artifactlinker}, a two-stage \emph{rank-then-verify} framework for automatic state-of-the-art discovery on Hugging Face. ArtifactLinker represents Hugging Face as a heterogeneous artifact graph with model, dataset, paper, and code-repository nodes, connected by edges such as \emph{fine-tune}, \emph{reference}, and \emph{evaluation} (with the measured metric, e.g. F1, as an edge attribute). A shared GNN encoder embeds every node from its neighbors; two heads read the same embedding jointly during training: a link head that estimates $P(\text{compatible})$ for a (model, dataset) pair, and an attribute head that estimates the expected evaluation score $\mathbb{E}[\text{score}]$ if that pair were evaluated. The two are trained together so that the link head, which sees negative (unobserved or incompatible) pairs, teaches the attribute head which pairs are viable to begin with -- a signal the attribute head could not learn from observed scores alone. The combined ranking score for an unobserved pair is $\text{rank} = \text{link} \times \text{attribute}$. ArtifactLinker's second stage lets an LLM coding agent reproduce the actual benchmark for only the top-ranked candidates, verifying predictions instead of trusting them outright.

This reframing -- model selection as link prediction on a joint artifact graph, following the same problem framing pioneered by TransferGraph's model-zoo graph learning \citep{li2024transfergraph} -- is precisely what our Retrieval Agent reuses as a tool: we load ArtifactLinker's trained checkpoint and query its link $\times$ attribute score for arbitrary (model, dataset) pairs, including a dataset we insert ourselves. What ArtifactLinker does not provide out of the box is (a) a way to place a genuinely new, unseen dataset into the graph, and (b) a mechanism to act on verification failures by retrieving again; both are the subject of Section~\ref{sec:approach}.
